\documentclass[11pt]{article}
\newcommand{\shorttitle}{Topic 0: Mathematical background}
\usepackage[T1]{fontenc}
\usepackage[utf8]{inputenc}
\usepackage{lmodern}
\usepackage[a4paper,margin=23mm,headheight=15pt]{geometry}
\usepackage{amsmath,amssymb,bm,graphicx,booktabs,longtable,array}
\usepackage{xcolor,microtype,enumitem,fancyhdr,fancyvrb,titlesec,tikz}
\usepackage[breakable]{tcolorbox}
\usepackage{hyperref}
\definecolor{navy}{HTML}{163653}
\definecolor{teal}{HTML}{007F86}
\definecolor{pale}{HTML}{EEF5F7}
\hypersetup{colorlinks=true,linkcolor=navy,urlcolor=teal,pdfauthor={Study notes based on the supplied teacher slides}}
\pagestyle{fancy}\fancyhf{}
\fancyhead[L]{\small\sffamily\color{navy}Artificial Intelligence}
\fancyhead[R]{\small\sffamily\color{navy}\shorttitle}
\fancyfoot[C]{\small\thepage}
\renewcommand{\headrulewidth}{0.3pt}
\titleformat{\section}{\Large\sffamily\bfseries\color{navy}}{\thesection}{0.7em}{}
\titleformat{\subsection}{\large\sffamily\bfseries\color{teal}}{\thesubsection}{0.6em}{}
\titleformat{\subsubsection}{\normalsize\sffamily\bfseries\color{navy}}{\thesubsubsection}{0.6em}{}
\setcounter{tocdepth}{2}\setcounter{secnumdepth}{2}
\setlength{\parindent}{0pt}\setlength{\parskip}{6pt plus 1pt}
\linespread{1.06}\setlist{itemsep=3pt,topsep=4pt,leftmargin=1.5em}
\setlength{\emergencystretch}{3em}
\newtcolorbox{keyidea}[1]{breakable,colback=pale,colframe=teal,boxrule=0.6pt,arc=1mm,title={#1},fonttitle=\sffamily\bfseries,coltitle=white,colbacktitle=teal}
\newtcolorbox{careful}[1]{breakable,colback=white,colframe=navy,boxrule=0.6pt,arc=1mm,title={#1},fonttitle=\sffamily\bfseries,coltitle=navy,colbacktitle=pale}
\newcommand{\sources}[1]{{\small\color{teal}\textit{Teacher slides: #1}}\par\nopagebreak[4]}
\newcommand{\newsection}[1]{\par\begingroup\skip0=\pagegoal\advance\skip0 by-\pagetotal\ifdim\skip0<12\baselineskip\newpage\fi\endgroup\section{#1}}
\newcolumntype{P}[1]{>{\raggedright\arraybackslash}p{#1}}
\newcommand{\E}{\operatorname{E}}
\newcommand{\Var}{\operatorname{Var}}
\newcommand{\Cov}{\operatorname{Cov}}
\newcommand{\R}{\mathbb{R}}
\newcommand{\argmax}{\operatorname*{arg\,max}}
\newcommand{\argmin}{\operatorname*{arg\,min}}
\newcommand{\relu}{\operatorname{ReLU}}
\newcommand{\codeinline}[1]{\texttt{\detokenize{#1}}}
\DefineVerbatimEnvironment{code}{Verbatim}{fontsize=\small,samepage=true,frame=single,rulecolor=\color{teal},framesep=2.5mm}
\newcommand{\cover}[3]{%
\begin{titlepage}\vspace*{12mm}
{\sffamily\color{teal}\large ARTIFICIAL INTELLIGENCE\par}\vspace{9mm}
{\sffamily\bfseries\color{navy}\Huge TOPIC #1\par}\vspace{5mm}
{\sffamily\color{navy}\LARGE #2\par}\vspace{11mm}
{\large Explanatory study notes\par}
{\large Theory, worked examples and revision questions\par}\vspace{12mm}
\begin{keyidea}{What you should be able to do}#3\end{keyidea}
\vfill
Based on the teacher PDFs supplied in \texttt{Lecture Slides}. The notes explain and reorganize that material. Worked calculations are study examples derived from the slide concepts, unless explicitly identified as a slide exercise. Clarifications of ambiguous or incorrect expressions are marked in context. References use \textbf{PDF page numbers}, counted from 1, rather than printed slide numbers.
\par\vspace{4mm}{\small English edition\enspace\textbullet\enspace 2 October 2026}
\end{titlepage}\setcounter{page}{2}}

\hypersetup{pdftitle={Topic 0 - Mathematical Background for AI}}
\begin{document}
\cover{0}{Mathematical background for AI}{Use conditioning, marginalization and Bayes' rule; calculate expectations and Gaussian predictions; connect entropy, KL divergence and cross-entropy to learning; understand how constraints become regularization.}
\tableofcontents\clearpage
\newsection{Why these tools appear throughout the course}
\sources{0.1. Math background, pp. 1--11.}
Machine learning models must handle incomplete information and noisy measurements. Probability represents uncertainty; information theory measures how well a model describes observations; optimization chooses parameters. These are connected: a probability model defines a likelihood, the negative log-likelihood becomes a loss, and an optimizer minimizes it.

Use $p(x)$ for a probability mass or density, depending on whether $X$ is discrete or continuous. Uppercase $X$ denotes a random variable; lowercase $x$ is one realized value. A vector is written in bold when helpful. $\E$ is expectation, $\Var$ is variance, and $\Sigma$ is a covariance matrix. All logarithms in these notes are natural unless a base is stated.
\begin{keyidea}{A chain of ideas to remember}
Probability model $\longrightarrow$ likelihood $\longrightarrow$ negative log-likelihood $\longrightarrow$ loss $\longrightarrow$ optimization. A prior adds assumptions about plausible parameters and can produce a regularization term.
\end{keyidea}
\newsection{Conditioning, joint probability and marginalization}
\sources{0.1. Math background, pp. 1--2.}
\subsection{Conditioning changes the information available}
For an event $B$ with positive probability,
\[
P(A\mid B)=\frac{P(A\cap B)}{P(B)},\qquad
P(A\cap B)=P(A\mid B)P(B).
\]
Conditioning restricts attention to outcomes where $B$ occurred. It is directional: $P(A\mid B)$ and $P(B\mid A)$ usually differ. Independence means $P(A\cap B)=P(A)P(B)$, so conditioning on $B$ does not change the probability of $A$.

For a sequence, the chain rule gives
\[
p(x_1,\ldots,x_n)=p(x_1)\prod_{i=2}^{n}p(x_i\mid x_1,\ldots,x_{i-1}).
\]
The chain rule always holds; simplifying each conditional to $p(x_i)$ requires independence. This distinction will matter for RL trajectories, where successive states are dependent.
\subsection{Worked slide exercise: drawing balls}
The slides give an urn containing three red balls and one blue ball. Draw without replacement. The chance of two reds is
\[
P(R_1,R_2)=P(R_1)P(R_2\mid R_1)=\frac34\frac23=\frac12.
\]
The second probability is $2/3$, because one red has already been removed. Using $3/4$ twice would incorrectly assume independence.

To find the chance that the second ball is red, sum over the possible first draws:
\begin{align*}
P(R_2)&=P(R_2\mid R_1)P(R_1)+P(R_2\mid B_1)P(B_1)\\
&=\frac23\frac34+1\frac14=\frac34.
\end{align*}
This is marginalization: remove a variable by summing over all its possible values. For a partition $B_1,\ldots,B_K$,
\[
P(A)=\sum_jP(A\mid B_j)P(B_j).
\]
For densities, replace the sum with an integral: $p(x)=\int p(x,y)\,dy$.
\newsection{Bayes' rule: reverse a conditional probability}
\sources{0.1. Math background, p. 2.}
\[
P(A\mid B)=\frac{P(B\mid A)P(A)}{P(B)}.
\]
In model language, the posterior is proportional to likelihood times prior:
\[
p(h\mid D)=\frac{p(D\mid h)p(h)}{\sum_jp(D\mid h_j)p(h_j)}.
\]
The \textbf{prior} describes uncertainty before these data; the \textbf{likelihood} describes how the assumed model generates these data; the \textbf{posterior} incorporates both. The denominator, called evidence, normalizes the result. Continuous parameters require integration instead of summation.
\subsection{Worked model-identification example}
Choose with equal probability between a fair coin $F$ and a biased coin $B$ with $P(H\mid B)=0.9$. Observe one head. Then
\[
P(B\mid H)=\frac{0.9(0.5)}{0.9(0.5)+0.5(0.5)}=\frac9{14}\approx0.643.
\]
One head provides evidence for the biased coin but does not identify it with certainty. After two independent heads conditional on the chosen coin,
\[
P(B\mid HH)=\frac{0.9^2}{0.9^2+0.5^2}\approx0.764.
\]
The two draws are conditionally independent given the coin. Their marginal dependence comes from sharing the same unknown coin.
\newsection{Random variables, expectation and variance}
\sources{0.1. Math background, pp. 2--4.}
\subsection{Mass, density and cumulative probability}
A discrete variable has probabilities $p(x)\geq0$ that sum to one. A continuous variable has a density $p(x)\geq0$ whose integral is one. For a continuous variable,
\[
P(a\leq X\leq b)=\int_a^b p(x)\,dx,\qquad
F(x)=P(X\leq x)=\int_{-\infty}^{x}p(t)\,dt.
\]
At differentiable points, $F'(x)=p(x)$. A density value is not a point probability: for a continuous variable $P(X=x)=0$, and a density can exceed one on a narrow interval.
\subsection{Expectation is a probability-weighted average}
\[
\E[g(X)]=\sum_xg(x)p(x)\quad\hbox{or}\quad
\E[g(X)]=\int g(x)p(x)\,dx.
\]
Linearity gives $\E[aX+bY]=a\E[X]+b\E[Y]$ without needing independence. In contrast, $\E[XY]=\E[X]\E[Y]$ generally needs an assumption such as independence.
\[
\mu=\E[X],\qquad \Var(X)=\E[(X-\mu)^2]=\E[X^2]-\mu^2.
\]
For two variables,
\[
\Var(aX+bY)=a^2\Var(X)+b^2\Var(Y)+2ab\Cov(X,Y).
\]
Independence removes the covariance term. Zero covariance alone does not generally imply independence.
\subsection{Bernoulli and indicator variables}
For $X\in\{0,1\}$ with success probability $\theta$,
\[
p(x)=\theta^x(1-\theta)^{1-x},\quad
\E[X]=\theta,\quad\Var(X)=\theta(1-\theta).
\]
An indicator $I_A$ is one on event $A$ and zero elsewhere, so $\E[I_A]=P(A)$. Logistic classification will use Bernoulli outputs, while nearest-neighbour voting can be written as a sum of class indicators.
\newsection{Gaussian distributions and linear transformations}
\sources{0.1. Math background, pp. 4--6.}
\subsection{Univariate normal}
\[
X\sim\mathcal N(\mu,\sigma^2),\qquad
p(x)=\frac{1}{\sqrt{2\pi\sigma^2}}
\exp\!\left[-\frac{(x-\mu)^2}{2\sigma^2}\right].
\]
The second argument is the \emph{variance}, not the standard deviation. Standardization gives $Z=(X-\mu)/\sigma\sim\mathcal N(0,1)$. Conversely, $X=\mu+\sigma Z$ has mean $\mu$ and variance $\sigma^2$.

Independent Gaussian variables have a Gaussian sum. Its mean is the sum of means and its variance is the sum of variances. More generally, a linear transformation of a jointly Gaussian vector is Gaussian.
\subsection{Multivariate normal}
For $\mathbf x\in\R^d$ and positive-definite covariance $\Sigma$,
\[
p(\mathbf x)=\frac{\exp[-\tfrac12(\mathbf x-\bm\mu)^T\Sigma^{-1}(\mathbf x-\bm\mu)]}
{(2\pi)^{d/2}|\Sigma|^{1/2}}.
\]
The quadratic form measures distance while accounting for scale and correlations. The log-density is
\[
\log p(\mathbf x)=-\frac d2\log(2\pi)-\frac12\log|\Sigma|
-\frac12(\mathbf x-\bm\mu)^T\Sigma^{-1}(\mathbf x-\bm\mu).
\]
For $Y=A\mathbf X+\mathbf b$, mean and covariance become $A\bm\mu+\mathbf b$ and $A\Sigma A^T$. If $\mathbf Z\sim\mathcal N(0,I)$ and $LL^T=\Sigma$, then $\bm\mu+L\mathbf Z\sim\mathcal N(\bm\mu,\Sigma)$. A Cholesky factor $L$ is a valid sampling factor, even though it is generally not the symmetric matrix square root.
\subsection{Conditioning a joint Gaussian}
Partition the mean and covariance into blocks for $X_a,X_b$. Then
\begin{align*}
X_b\mid X_a=x_a&\sim\mathcal N(\mu_{b\mid a},\Sigma_{b\mid a}),\\
\mu_{b\mid a}&=\mu_b+\Sigma_{ba}\Sigma_{aa}^{-1}(x_a-\mu_a),\\
\Sigma_{b\mid a}&=\Sigma_{bb}-\Sigma_{ba}\Sigma_{aa}^{-1}\Sigma_{ab}.
\end{align*}
The conditional mean adjusts according to correlation; the covariance decreases because an observation adds information. This identity is the basis of Gaussian-process prediction.
\begin{careful}{Clarifications of the Gaussian summary on pp. 5--6}
The conditional covariance uses a \textbf{minus} sign. A weighted mixture of Gaussian densities is generally not Gaussian. Gaussian marginals alone do not guarantee a jointly Gaussian distribution. A linear combination of correlated Gaussian blocks must include cross-covariance terms. Also, $\mu+\Sigma^{1/2}Z$ has covariance $\Sigma$, rather than $\Sigma^2$. These distinctions prevent errors when using the compact slide summary.
\end{careful}
\subsection{Worked conditioning example}
Let $(X,Y)$ have zero means, unit variances and covariance $0.8$. For $X=1$,
\[
Y\mid X=1\sim\mathcal N(0.8,1-0.8^2)=\mathcal N(0.8,0.36).
\]
The conditional standard deviation is $0.6$. The prediction is uncertain even after $X$ is known, because correlation is not perfect.
\newsection{Entropy, cross-entropy and KL divergence}
\sources{0.1. Math background, pp. 6--9.}
\subsection{Surprise and entropy}
The surprise of an outcome with probability $p(x)$ is $-\log p(x)$. Entropy is expected surprise:
\[
H(p)=-\sum_xp(x)\log p(x).
\]
With $\log_2$, units are bits; with natural logs, they are \textbf{nats} (the handout calls them ``nits''). Use the convention $0\log0=0$.

For a Bernoulli variable,
\[
H(\theta)=-\theta\log\theta-(1-\theta)\log(1-\theta).
\]
A fair coin has entropy $\log2\approx0.693$ nats, or one bit. A deterministic coin has entropy zero. Differential entropy for continuous variables replaces the sum by an integral, but can be negative and depends on the measurement scale.
\subsection{Cross-entropy evaluates a model against data}
If data follow $p$ and a model assigns probabilities $q$,
\[
H(p,q)=-\sum_xp(x)\log q(x),\qquad
\mathrm{KL}(p\|q)=\sum_xp(x)\log\frac{p(x)}{q(x)}.
\]
Therefore,
\[
\boxed{H(p,q)=H(p)+\mathrm{KL}(p\|q)}.
\]
KL is nonnegative and is zero when distributions agree, apart from null sets. It is not symmetric and does not satisfy the requirements of a distance metric. If $q$ assigns zero probability to an event with positive $p$, the divergence is infinite.
\subsection{Worked slide exercise: two coins}
For $P=(0.5,0.5)$ and $Q=(0.9,0.1)$,
\begin{align*}
\mathrm{KL}(P\|Q)&=0.5\log(0.5/0.9)+0.5\log(0.5/0.1)\approx0.511,\\
\mathrm{KL}(Q\|P)&=0.9\log(0.9/0.5)+0.1\log(0.1/0.5)\approx0.368.
\end{align*}
The larger first value reflects the cost of using a model that treats tails as rare when tails actually occur half the time.
\subsection{Information gain: a careful distinction}
The handout links posterior-versus-prior KL to entropy reduction. For a \emph{particular} observed $y$, define the realized Bayesian information gain as $\mathrm{KL}(p(X\mid y)\|p(X))$. It is not generally equal to $H(X)-H(X\mid Y=y)$.

The equality holds \emph{after averaging over observations}:
\[
I(X;Y)=\E_Y[\mathrm{KL}(p(X\mid Y)\|p(X))]
=H(X)-H(X\mid Y).
\]
This expected entropy reduction is the information-gain idea used in decision-tree splitting. A particular surprising observation can increase posterior entropy; the expected reduction is nonnegative.
\newsection{Likelihood, estimation and optimization}
\sources{0.1. Math background, pp. 8--11.}
\subsection{Maximum likelihood and cross-entropy}
For independent observations conditional on parameters $w$,
\[
L(w)=\prod_{i=1}^n p(x_i\mid w),\qquad
\hat w_{\mathrm{ML}}=\argmax_w\sum_i\log p(x_i\mid w).
\]
Taking logs preserves the maximizer and turns products into sums. Dividing by $n$ does not alter the maximizer either. The negative average log-likelihood is the empirical cross-entropy between the observed distribution and the model. Its population counterpart differs from KL only by the entropy of the true distribution, which is independent of $w$.

Maximum a posteriori estimation instead maximizes $\log L(w)+\log p(w)$. Full Bayesian inference retains $p(w\mid D)$ rather than only its maximizing point.
\begin{careful}{A qualification to the MLE statement on p. 9}
MLE need not be unbiased in a finite sample. Consistency and asymptotic normality require regularity and correct-model conditions. With per-observation Fisher information $I(w_*)$, the usual statement is $\sqrt n(\hat w-w_*)\Rightarrow\mathcal N(0,I(w_*)^{-1})$. Identifiability and sample size matter.
\end{careful}
\subsection{A norm constraint becomes ridge regularization}
Start from
\[
\min_w f(w)\quad\hbox{subject to}\quad\|w\|_2^2\leq C.
\]
The Lagrangian is $\mathcal L(w,\lambda)=f(w)+\lambda(\|w\|_2^2-C)$ with $\lambda\geq0$. For fixed $\lambda$, the term $-\lambda C$ does not depend on $w$, giving a penalized objective $f(w)+\lambda\|w\|_2^2$.

At a differentiable optimum, stationarity is $\nabla f(w)+2\lambda w=0$. The full constraint conditions also require feasibility and complementary slackness: $\lambda(\|w\|^2-C)=0$. An inactive constraint can have $\lambda=0$; the tangency picture describes the active case. The correspondence between $C$ and $\lambda$ is especially useful in convex regression.
\newsection{Revision sheet and self-test}
\begin{longtable}{@{}P{.30\textwidth}P{.61\textwidth}@{}}
\toprule Tool & What to ask yourself\\\midrule\endhead
Conditioning & What information is fixed, and is the conditioning event possible?\\
Marginalization & Which variable must be removed by a sum or integral?\\
Bayes & What are prior, likelihood and normalizing evidence?\\
Expectation & What outcomes are being averaged, under which distribution?\\
Gaussian conditioning & Are covariance blocks ordered correctly, with subtraction?\\
Cross-entropy & Which distribution generates data and which predicts them?\\
Regularization & What assumption does the penalty impose?\\\bottomrule
\end{longtable}
\subsection{Questions}
\begin{enumerate}
\item In the urn, what is $P(R_1\mid R_2)$?
\item Does linearity of expectation need independent variables?
\item If $X\sim\mathcal N(2,9)$, how do you standardize it?
\item Can a mixture of two Gaussians have two peaks?
\item What is cross-entropy for a correct class assigned probability $0.1$?
\item Why can cross-entropy be minimized without knowing $H(p)$?
\item Does a smaller norm budget normally correspond to stronger regularization?
\item Must one particular observation decrease posterior entropy?
\end{enumerate}
\subsection{Answers}
\begin{enumerate}
\item $P(R_1,R_2)/P(R_2)=(1/2)/(3/4)=2/3$.
\item No; linearity always holds when the expectations exist.
\item $(X-2)/3$; divide by SD, not variance.
\item Yes. A Gaussian mixture need not be Gaussian.
\item $-\log0.1\approx2.303$ nats.
\item $H(p)$ is constant with respect to the model parameters.
\item Yes, in the usual convex constraint--penalty correspondence.
\item No. Entropy reduction is nonnegative in expectation.
\end{enumerate}
\newsection{Source guide}
\texttt{0.1. Math background.pdf}: probability (pp. 1--6), information theory and MLE (pp. 6--9), constrained optimization (pp. 9--11). These notes derive the handout's urn and coin exercises and clarify its compact Gaussian and information-gain summaries. No external reading is required for the explanations here.
\end{document}
